\documentclass[xcolor=dvipsnames]{beamer}
\usetheme{nelle}
\usepackage{natbib}                 % Fancy bibliography.
\usepackage{url}                    % Allow printing of URLs.
\usepackage{outlines}
\usepackage{enumitem}
\usepackage{multicol}
\usepackage{dsfont}
\usepackage{amsmath}
\usepackage{epstopdf}
\usepackage{listings}
\usepackage{color}
\setbeamerfont{caption}{size=\scriptsize}
\setbeamertemplate{navigation symbols}{}
\setbeamertemplate{footline}[frame number]{}

\def\newblock{\hskip .11em plus .33em minus .07em}
\usepackage{color}

\definecolor{mygreen}{rgb}{0,0.6,0}
\definecolor{mygray}{rgb}{0.5,0.5,0.5}
\definecolor{mymauve}{rgb}{0.58,0,0.82}

\lstset{ %
  backgroundcolor=\color{white},   % choose the background color; you must add \usepackage{color} or \usepackage{xcolor}; should come as last argument
  basicstyle=\scriptsize,        % the size of the fonts that are used for the code
  breakatwhitespace=false,         % sets if automatic breaks should only happen at whitespace
  breaklines=true,                 % sets automatic line breaking
  captionpos=b,                    % sets the caption-position to bottom
  commentstyle=\color{mygreen},    % comment style
  deletekeywords={...},            % if you want to delete keywords from the given language
  escapeinside={\%*}{*)},          % if you want to add LaTeX within your code
  extendedchars=true,              % lets you use non-ASCII characters; for 8-bits encodings only, does not work with UTF-8
  keepspaces=true,                 % keeps spaces in text, useful for keeping indentation of code (possibly needs columns=flexible)
  keywordstyle=\color{blue},       % keyword style
  language=Python,                 % the language of the code
  morekeywords={*,...},            % if you want to add more keywords to the set
  rulecolor=\color{black},         % if not set, the frame-color may be changed on line-breaks within not-black text (e.g. comments (green here))
  showspaces=false,                % show spaces everywhere adding particular underscores; it overrides 'showstringspaces'
  showstringspaces=false,          % underline spaces within strings only
  showtabs=false,                  % show tabs within strings adding particular underscores
  stringstyle=\color{mymauve},     % string literal style
  tabsize=2,                     % sets default tabsize to 2 spaces
  title=\lstname                   % show the filename of files included with \lstinputlisting; also try caption instead of title
}

\newcommand{\todo}[1]{\textbf{[TODO: #1]}}
\newcommand{\fixme}[1]{\textbf{[FIXME: #1]}}


\title{\textbf{The Scientific Python Toolstack for biomedical research}}

\author[Varoquaux Nelle]{
Nelle Varoquaux}


\date{December, 3rd}
\institute{Mines ParisTech, Institut Curie, INSERM}
\begin{document}
\begin{frame}[t, noframenumbering]
  \maketitle

\end{frame}

\setcounter{framenumber}{0}

\section{A bit about me \& OpenSource}
\begin{frame}
\frametitle{A bit about me\dots}
\begin{columns}
\begin{column}{0.45\linewidth}
\centering
\includegraphics[width=0.5\linewidth]{images/me.png}
\end{column}
\begin{column}{0.6\linewidth}
\begin{overlayarea}{\linewidth}{6cm}
\only<1>{
\vspace{3em}
\begin{itemize}[label={$\bullet$}]
\item Engineering degree, specialized in Software Engineering
\item Masters in Machine Learning
\item PhD in computational biology
\item Now a postdoc in the Statistics Department at UC Berkeley and BIDS
\end{itemize}}

\only<2>{
\vspace{1em}
\begin{itemize}[label={$\bullet$}]
\item Contributed to
\begin{itemize}[label={$\bullet$}]
\small
\item scikit-learn
\item Matplotlib
\item MarkUs
\item OpenOffice\dots
\end{itemize}
\item Helped organize:
\begin{itemize}[label={$\bullet$}]
\small
\item Euroscipy 2010, 2011, 2012, 2013, \dots
\item Scipy 2015, 2016
\item PyconFr 2011, 2012, 2013, 2014, \dots
\end{itemize}
\item Chaired the French Python association between 2012 \& 2015
\end{itemize}
}
\end{overlayarea}
\end{column}
\end{columns}
{\footnotesize
\vspace{3em}
\begin{center}
Twitter: @nvaroqua \hspace{3em} Github: @NelleV
\end{center}
}
\end{frame}


\begin{frame}

\begin{columns}
\begin{column}{0.6\linewidth}
{\Large \bf Opensource, and the Python community \& ecosystem was a big part of my success in academia}
\end{column}
\begin{column}{0.4\linewidth}
\begin{flushright}
\includegraphics[width=0.9\linewidth]{figures/obelix.png}
\end{flushright}
\end{column}
\end{columns}
\end{frame}

\begin{frame}
\frametitle{As a matter of fact\dots}
\begin{center}
\includegraphics[width=0.2\linewidth]{images/fernando.png}
\hspace{2em}
\includegraphics[width=0.2\linewidth]{images/stefan.jpg}
\hspace{2em}
\includegraphics[width=0.2\linewidth]{images/matthias.jpg}
\end{center}

\end{frame}

\begin{frame}
\frametitle{A very diverse community}

{\Large ``
A very diverse community [..]: \\
I am talking to you about things you don't understand (my science) \\
and things I don't understand (web dev)
''}
\begin{flushright}
\em Gael Varoquaux \\
\texttt{scikit-learn} lead
\end{flushright}
\end{frame}


\section{The scientific Python toolstack for biomedical research}
\begin{frame}
\Large{ \bf
The scientific Python toolstack for biomedical research}

\end{frame}

\begin{frame}
\frametitle{The DNA sequence holds the genetic information}
\begin{columns}
\begin{column}{0.45\linewidth}
\includegraphics[width=\linewidth]{images/bio_101.png}
\end{column}
\begin{column}{0.45\linewidth}
\begin{itemize}[label={$\bullet$}]
\item Beings are composed of one or several cells.
\item Cells contain the genetic code of an individual on the DNA sequence.
\item This DNA sequence is roughly identical between organisms, and between
cells of an organism.
\end{itemize}
\end{column}
\end{columns}
\end{frame}

\begin{frame}
\frametitle{Several genome sizes}

\begin{table}
\footnotesize
\begin{tabular}{l | l | l | c}
{\bf Organisms} &{\bf Chromosomes} & {\bf Size} & {\bf Ratio} \\
\hline
{\bf Bacteria} & 1 & 400,000 to 10,000,000 & $\leq 1$ \\
{\color{Blue} \bf In search of lost time} & 7 & 9,609,000 & 1\\
{\bf Yeast} & 12 & 14,000,000 & 1.5 \\
{\bf Fly} & 4 & 300,000,000 & 30 \\
{\bf Human} & 46 & 6,000,000,000 & 600 \\
{\bf Tulip} & 60 & 30,000,000,000 & 3,000 \\
\end{tabular}
\end{table}
\end{frame}

\begin{frame}
\frametitle{As sequencing costs diminish\dots}
\begin{figure}
\includegraphics[width=0.9\linewidth]{images/dna_sequences_growth.PNG}
\end{figure}
\end{frame}


\begin{frame}
\frametitle{How can we analyse such large quantity of data?}
\begin{figure}
\includegraphics[width=0.9\linewidth]{images/dna_sequences_growth.PNG}
\end{figure}
\end{frame}

\begin{frame}
\frametitle{Can we use the Scientific Python Toolstack?}

\begin{center}
\includegraphics[width=0.2\linewidth]{images/pythonlogo.jpg} \hspace{3em}
\includegraphics[width=0.2\linewidth]{images/numpy.jpg}\hspace{3em}
\includegraphics[width=0.2\linewidth]{images/scipy.jpg}
\end{center}

\vspace{1.5em}

\begin{center}
\hspace{3em}
\includegraphics[width=0.2\linewidth]{images/scikit-learn.png}\hspace{3em}
\includegraphics[width=0.2\linewidth]{images/ipython.jpg}\hspace{3em}
\includegraphics[width=0.2\linewidth]{images/matplotlib.png}
\end{center}
\end{frame}

\begin{frame}
\frametitle{\texttt{numpy}: fast numerical arrays computations}
\begin{center}
\includegraphics[width=0.2\linewidth]{images/numpy.jpg}
\end{center}
\only<1>{
\lstinputlisting{scripts/numpy.py}}
\only<2>{
\begin{center}
\includegraphics[width=0.6\linewidth]{images/numpy_three_concepts.png}
\\
numpy array = $\bullet$ memory adress \hspace{1em} $\bullet$ data type \hspace{1em} $\bullet$ shape
\hspace{1em} $\bullet$ stride \\
\vspace{2em}
{\bf \footnotesize Fast in-memory operations, thanks to numerical libraries such as blas}
\end{center}
}
\end{frame}


\begin{frame}
\frametitle{\texttt{scipy}: a scientific computing package for Python}
\begin{overlayarea}{12cm}{3cm}
\begin{center}
\includegraphics[width=0.2\linewidth]{images/scipy.jpg}
\end{center}
\end{overlayarea}
\begin{overlayarea}{12cm}{6cm}
\only<1>{
\begin{columns}
\begin{column}{0.3\linewidth}
\begin{itemize}[label={$\bullet$}]
\item cluster
\item fftpack
\item integrate
\item interpolate
\item io
\end{itemize}
\end{column}
\begin{column}{0.3\linewidth}
\begin{itemize}[label={$\bullet$}]
\item linalg
\item misc
\item ndimage
\item odr
\item optimize
\end{itemize}
\end{column}
\begin{column}{0.3\linewidth}
\begin{itemize}[label={$\bullet$}]
\item signal
\item sparse
\item spatial
\item special
\item stats
\end{itemize}
\end{column}
\end{columns}}
\only<2>{
\begin{columns}
\begin{column}{0.3\linewidth}
\begin{itemize}[label={$\bullet$}]
\item {\bf cluster}
\item fftpack
\item integrate
\item interpolate
\item io
\end{itemize}
\end{column}
\begin{column}{0.3\linewidth}
\begin{itemize}[label={$\bullet$}]
\item {\bf linalg}
\item misc
\item ndimage
\item odr
\item {\bf optimize}
\end{itemize}
\end{column}
\begin{column}{0.3\linewidth}
\begin{itemize}[label={$\bullet$}]
\item signal
\item {\bf sparse}
\item spatial
\item special
\item {\bf stats}
\end{itemize}
\end{column}
\end{columns}
}
\only<3>{
\lstinputlisting{scripts/scipy_optimize.py}
}
\end{overlayarea}
\end{frame}

\begin{frame}
\frametitle{\texttt{IPython}: an interactive python interpreter}
\begin{center}
\includegraphics[width=0.7\linewidth]{images/ipython_013.png}
\end{center}
\end{frame}

\begin{frame}
\frametitle{\texttt{Matplotlib}: plotting library in Python}
\begin{center}
\includegraphics[width=0.5\linewidth]{images/sphx_glr_anatomy_001.png}
\end{center}
\end{frame}

\begin{frame}
\frametitle{\texttt{scikit-learn}: a machine learning toolkit in Python}
\begin{center}
\only<1>{
\includegraphics[width=0.5\linewidth]{images/sklearn_regression.png}
}
\only<2>{
\lstinputlisting{scripts/sklearn_regression.py}
}
\end{center}
\end{frame}



\begin{frame}
\frametitle{}
{\Large \bf A couple of examples of projects accomplished in Python}
\end{frame}

\begin{frame}
\frametitle{``Cell cognition''}
\begin{center}
\includegraphics[width=0.6\linewidth]{images/cell_cognition.png}
\end{center}
\end{frame}

\begin{frame}
\frametitle{DREAM challenge: predicting the toxicity of drugs}
\begin{center}
\includegraphics[width=0.7\linewidth]{images/toxicogenetic_challenge.png}
\end{center}
\end{frame}

\begin{frame}
\frametitle{What are you thinking of?}
\begin{center}
\includegraphics[width=0.5\linewidth]{images/brain_activity.png}
\end{center}
\end{frame}


\begin{frame}
\frametitle{My personal interest: the 3D structure of the genome}
\begin{columns}
\begin{column}{0.45\linewidth}
\includegraphics[width=\linewidth]{images/bio_101.png}
\end{column}
\begin{column}{0.45\linewidth}
\begin{itemize}[label={$\bullet$}]
\item On the average, a single human chromosome consists of DNA molecule that is about 5cm long
\item A nucleus is on average 6$\mu$m of diameter.
\item {\bf \color{Blue} How does DNA fit into the nucleus?}
\end{itemize}
\end{column}
\end{columns}
\end{frame}

\begin{frame}
\frametitle{The 3D structure of the genome is thought to play an important
role in many biological processes}
\vspace{-0.6em}
\begin{figure}
\begin{center}
\includegraphics[width=0.7\linewidth]{figures/yeasts_genome_architecture.jpg}
\end{center}
\caption{\textbf{The genome of \textit{S. cerevisiae} is highly organized}
         \citep{zimmer:principles}}
\end{figure}
\end{frame}

% 3. A contact count matrix
\begin{frame}
\frametitle{The contact count matrix recapitulates the hallmarks of genome
architecture}
\vspace{-1em}
\begin{figure}
\includegraphics[width=0.55\textwidth]{figures/yeast_counts.pdf}
\caption{\textbf{Contact counts for the first 5 chromosomes of \textit{S.
cerevisiae}}}
\end{figure}
\end{frame}

\begin{frame}
\frametitle{The computational challenges of Hi-C data}
\begin{itemize}[label={$\bullet$}]
\item The size of the organisms studied is growing.
\item The resolution is getting higher.
\item A 1~kb human map is 3,000,000 by 3,000,000.
\item In a naive format, {\bf this corresponds to 1.1~Tb of data}
\end{itemize}
\end{frame}


\section{ICE: how good data structure beats parallelization}
\begin{frame}
\Large{ \bf
\texttt{Iced}: Normalization of Hi-C data}

\begin{flushright}
\vspace{1em}
\small
\textit{How appropriate use of data structures beats C++ parallized code.}
\end{flushright}
\end{frame}

\begin{frame}
\frametitle{Contact-maps normalization: the ICE method}

{\bf \color{Blue} Sequencing data is full of biases}

\begin{itemize}[label={$\bullet$}]
\item Sequencing biases
\begin{itemize}
\item ("CG" riched regions are sequenced less than "AT" rich region)
\end{itemize}
\item The protocol induces more biases.
\end{itemize}
\vspace{3em}

\begin{center}
{\bf \color{Blue} We need to correct for these biases}
\end{center}

\end{frame}


\begin{frame}
\frametitle{Iterative Correction of Contact maps}
{\bf \color{Blue} Two hypotheses}

Denote by $C$ the raw contact count matrix and by $N$ the normalized one.

\begin{itemize}[label={$\bullet$}]
\item Each row and column interacts as much as any over row and column.
$$
\sum_{i} N_{ij} = K
$$
\item The bias of each entry of the matrix can be decomposed as the product of
a row and column-specific bias
$$
\beta_{ij} = \beta_i \beta_j
$$
\end{itemize}
\end{frame}

\begin{frame}
\frametitle{Formulating an optimization problem}
\begin{equation*}
\renewcommand{\arraystretch}{2}
\begin{array}{ccll}
\text{Find} & & \beta \\
\text{Such that} & & \underset{i}{\sum} \quad \beta_i \beta_j c_{ij} = K
\end{array}
\end{equation*}

\vspace{1em}
{\bf \color{Blue} Matrix-balancing problems} \\
\begin{itemize}[label={$\bullet$}]
\item This boils down to solving a {\bf matrix-balancing} problem.
\item These belong to an extremely widely studied class of problems.
\end{itemize}
\end{frame}

\begin{frame}
\frametitle{The Sinkhorn-Knopp algorithm}
{\bf \color{Blue} A very simple algorithm exists to solve this issue.}
\lstinputlisting{scripts/algorithm}
\end{frame}

\begin{frame}
\frametitle{A simple implementation}
\lstinputlisting{scripts/ice.py}
\end{frame}

\begin{frame}
\frametitle{A problem with the simple implementation}
\begin{center}
\begin{centering}
{\bf \color{Red} \Large This will not scale!}
\end{centering}
\end{center}
\end{frame}

\begin{frame}
\frametitle{Using scipy.sparse matrices !}
{\bf \color{Blue} Property of our data}
\begin{itemize}[label={$\bullet$}]
\item Our data is sparse (most of it is 0):
\item Our data is a symmetric matrix;
\item \texttt{scipy.sparse} can store sparse matrix and provide ``numpy-like'' interface.
\end{itemize}

\begin{columns}
\begin{column}{0.5\linewidth}
\end{column}
\begin{column}{0.5\linewidth}
\lstinputlisting{scripts/coo_data}
\end{column}
\end{columns}
\end{frame}

\begin{frame}
\frametitle{Using Sparse Compressed Row format}
{\bf \color{Blue} The CSR format}

\begin{itemize}[label={$\bullet$}]
\item \texttt{indices} is array of column indices
\item \texttt{data} is array of corresponding nonzero values
\item \texttt{indptr} points to row starts in indices and data
\end{itemize}
\end{frame}

\begin{frame}
\frametitle{Using Sparse Compressed Row format}
\begin{columns}
\begin{column}{0.5\linewidth}
\lstinputlisting{scripts/dense_format}
\end{column}
\begin{column}{0.5\linewidth}
\begin{itemize}[label={$\bullet$}]
\item indices: {\small \texttt{[2, 1, 3, 2, 4, 0, 2]}}
\item data: {\small \texttt{[ 1.015,  0.838,  1.724,  1.724, 1.787,  0.574,
1.153]}}
\item indptr: {\small \texttt{[0, 1, 3, 4, 5, 7]}}
\end{itemize}
\end{column}

\end{columns}
\end{frame}

\begin{frame}
\frametitle{More on the Sparse Compressed Row format}
\begin{itemize}[label={$\bullet$}]
\item length is \texttt{n\_row + 1}, last item = number of values = length of both indices and data
\item nonzero values of the i-th row are \texttt{data[indptr[i]:indptr[i+1]]} with column indices \texttt{indices[indptr[i]:indptr[i+1]]}
\item item $(i, j)$ can be accessed as \texttt{data[indptr[i]+k]}, where $k$ is position of $j$ in \texttt{indices[indptr[i]:indptr[i+1]]}
\end{itemize}
\end{frame}


\begin{frame}
\frametitle{In short\dots}
{\Large \bf \color{Blue} It's roughly the same as COO for sums over row.}\\
\vspace{1em}
{\Large \bf \color{Blue} But sometimes we need to loop over
elements to perform operations}\\
\vspace{1em}
{\Large \bf \color{Blue} \flushleft And looping is slow\dots}
\end{frame}

\begin{frame}
\frametitle{Cython to the rescue}
\only<1>{
\begin{center}
\includegraphics[width=0.5\linewidth]{images/cython.png}
\end{center}
}
\only<2>{

\lstinputlisting{scripts/cython_normalization.pyx}
}
\end{frame}


\begin{frame}
\frametitle{How does our Cython/Python implementation perform?}

{HiCorrector is a C++ implementation of the same algorithm, parallelized but using dense format.}

\vspace{2em}
\begin{table}
\scriptsize
\begin{tabular}{l | c | c | c | c }
  & {\bf Iced {\tiny (dense)}} & {\bf Iced {\tiny (sparse)}} & {\bf HiCorrector {\tiny (1 CPU)}} & {\bf HiCorrector {\tiny (8 CPUs)}} \\
\hline
IMR90 1~Mbp & 00:00:12 & 00:00:25 & 00:00:25 & 00:00:06 \\
IMR90 500~kbp & 00:00:40 & 00:01:30 & 00:02:15 & 00:00:22 \\
IMR90 150~kbp & - & 00:04:28 & 00:13:21 & 00:03:10 \\
IMR90 40~kbp & - & 00:07:19 & 02:35:34 & 00:35:43 \\
IMR90 20~kbp & - & 00:08:36 & 12:57:17 & 02:34:05 \\
\end{tabular}
\end{table}
\end{frame}

\begin{frame}
\frametitle{Conclusion}

{\Large \em Using better data structure can lead to massive performance gain.}

\vspace{1em}

\begin{itemize}[label={$\bullet$}]
\item Always use a programming language you are familiar with.
\item If you are not familiar with any programming language, use an ``easy''
programming language (R, Python, etc).
\item In doubt, benchmark to understand where the code is slow.
\item {\bf I was able to quickly implement a fast, competitive implementation
of the algorithm thanks to the Scientific Python toolstack}
\end{itemize}
\end{frame}

\begin{frame}
\frametitle{Acknowledgments}
\fboxsep=0pt
\noindent
\begin{minipage}[t]{0.48\linewidth}
\textbf{Mines ParisTech - CBIO} \\
Jean-Philippe \textsc{Vert} \\
Thomas \textsc{Walter} \\

\textbf{UW - Noble lab} \\
William S. \textsc{Noble} \\
Ferhat \textsc{Ay} \\

\textbf{Institut Curie - U900} \\
Emmanuel \textsc{Barillot} \\
Nicolas \textsc{Servant} \\
Chong-Jian \textsc{Chen} \\
Eric \textsc{Viara} \\

\textbf{UMass} \\
Job \textsc{Dekker} \\
Bryan \textsc{Lajoie} \\

\end{minipage}
\hfill%
\begin{minipage}[t]{0.48\linewidth}

\textbf{UC Riverside} \\
Karine \textsc{Le Roch} \\
Sebastiaan \textsc{Bol} \\
Evelien \textsc{Bunnik} \\
Jacques \textsc{Prudhomme} \\

\textbf{Institut Curie - UMR168} \\
Edith \textsc{Heard} \\

\textbf{UW - Dunham lab} \\
Maitreya \textsc{Dunham} \\
Ivan \textsc{Liachko} \\

\textbf{UW - Shendure lab} \\
Jay \textsc{Shendure} \\
Josh \textsc{Burton} \\

\end{minipage}

\end{frame}

\begin{frame}[allowframebreaks,noframenumbering]
  \frametitle{References}
  \bibliographystyle{plainnat}
  \bibliography{refs}
\end{frame}

\end{document}

